\documentclass[12pt,a4paper,oneside]{report}

% -----------------------------------------------------------------------------
% Configurarea documentului
% -----------------------------------------------------------------------------
\usepackage[utf8]{inputenc}
\usepackage[T1]{fontenc}
\usepackage[romanian]{babel}
\usepackage{lmodern}
\usepackage{geometry}
\usepackage{microtype}
\usepackage{amsmath,amssymb}
\usepackage{siunitx}
\usepackage{booktabs}
\usepackage{array}
\usepackage{longtable}
\usepackage{enumitem}
\usepackage{xcolor}
\usepackage{listings}
\usepackage{hyperref}

\geometry{left=30mm,right=25mm,top=25mm,bottom=25mm}
\sisetup{output-decimal-marker={,}}
\hypersetup{
  colorlinks=true,
  linkcolor=blue!50!black,
  urlcolor=blue!60!black,
  pdftitle={Control PID cascadat -- proiectare, reglare și validare},
  pdfauthor={Alin Seicarin}
}

\lstset{
  basicstyle=\ttfamily\small,
  frame=single,
  breaklines=true,
  columns=fullflexible,
  showstringspaces=false
}

\newcommand{\sat}{\operatorname{sat}}
\newcommand{\wrap}{\operatorname{wrap}_{[-\pi,\pi]}}

\title{%
  \textbf{Compararea strategiilor de control al traiectoriei pentru un robot
  diferențial simulat în Gazebo}\\[1.5em]
  \large Proiectarea, reglarea și validarea regulatorului PID cascadat
}
\author{Alin Seicarin}
\date{Iulie 2026}

\begin{document}

\maketitle
\tableofcontents
\clearpage

% =============================================================================
\chapter{Introducere și obiective}
% =============================================================================

Tema proiectului este compararea strategiilor de control aplicate unui robot
mobil cu acționare diferențială, simulat în Gazebo. Variantele principale care
vor fi analizate sunt regulatorul PID, regulatorul liniar pătratic (LQR) și
controlul predictiv bazat pe model (MPC). În etapa descrisă în acest document au
fost dezvoltate infrastructura experimentală, un regulator PID preliminar bazat
pe urmărirea unui punct anticipat și regulatorul PID cascadat propus pentru
comparația principală.

Obiectivele urmărite sunt:

\begin{enumerate}[label=\arabic*.]
  \item menținerea robotului în vecinătatea unei căi plane predefinite;
  \item folosirea acelorași referințe geometrice, limite de comandă și criterii
        de evaluare pentru toate regulatoarele;
  \item măsurarea erorii laterale, erorii de orientare, duratei și activității
        normalizate a comenzilor;
  \item evaluarea pe trasee cu proprietăți geometrice diferite și în prezența
        unei perturbații reproductibile a comenzii unghiulare;
  \item păstrarea experimentelor reproductibile prin fișiere de configurare și
        rulări automatizate.
\end{enumerate}

\section{Urmărirea căii și urmărirea traiectoriei}

În forma curentă, problema implementată este una de \emph{urmărire geometrică a
căii} (path following). Calea stabilește poziția și orientarea dorite în funcție
de coordonata curbilinie $s$, iar managerul de referință furnizează și o viteză
nominală $v_{\mathrm{ref}}(s)$. Nu este impusă o lege temporală strictă de
forma $(x_{\mathrm{ref}}(t),y_{\mathrm{ref}}(t))$.

Această alegere fixează domeniul lucrării deoarece permite compararea
regulatoarelor în aceleași coordonate de eroare, fără ca o abatere temporară să
producă automat o eroare longitudinală mare. Termenul „urmărirea traiectoriei”
din titlul temei este folosit în sensul larg al urmăririi unei căi geometrice;
nu se adaugă ulterior o parametrizare temporală.

% =============================================================================
\chapter{Platforma de simulare}
% =============================================================================

\section{Robotul diferențial}

Modelul conține un șasiu rigid, două roți motoare și un punct de sprijin frontal.
Raza fiecărei roți motoare este $r=\SI{0.1}{m}$, iar distanța dintre centrele
roților este $L=\SI{0.35}{m}$. Masa șasiului este de \SI{3}{kg}, iar centrul său
de masă este deplasat cu \SI{0.05}{m} în fața centrului geometric. Această
asimetrie elimină caracterul ideal perfect simetric al modelului.

Pentru vitezele unghiulare ale roților $\omega_R$ și $\omega_L$, viteza liniară
și viteza unghiulară ale corpului sunt:

\begin{align}
  v &= \frac{r}{2}\left(\omega_R+\omega_L\right), \\
  \omega &= \frac{r}{L}\left(\omega_R-\omega_L\right).
\end{align}

Modelul cinematic plan utilizat pentru interpretarea mișcării este:

\begin{align}
  \dot{x} &= v\cos\psi, \\
  \dot{y} &= v\sin\psi, \\
  \dot{\psi} &= \omega,
\end{align}

unde $(x,y)$ reprezintă poziția cadrului robotului, iar $\psi$ este unghiul de
orientare. Constrângerea neholonomă exclude deplasarea laterală independentă.

\section{Gazebo, ROS~2 și estimarea stării}

Simularea folosește ROS~2 Humble și Gazebo Classic. Pluginul de acționare
diferențială primește comenzi de tip \texttt{geometry\_msgs/Twist} pe tema
\texttt{/cmd\_vel} și publică pe \texttt{/odom} o odometrie obținută prin
integrarea vitezelor articulațiilor roților. Sursa odometriei a fost configurată
explicit în modul \emph{encoder}; prin urmare, poziția publicată nu mai este
citită direct din starea exactă a lumii Gazebo. Parametrii geometrici ai
pluginului sunt identici cu geometria URDF: diametrul roții este \SI{0.2}{m},
iar separarea roților este \SI{0.35}{m}.

Un senzor IMU simulat publică pe \texttt{/demo/imu}. Nodul
\texttt{robot\_localization/ekf\_node} combină viteza longitudinală, condiția
neholonomă $v_y=0$ și viteza unghiulară furnizate de odometria roților cu viteza
unghiulară și accelerația longitudinală măsurate de IMU. Poziția și viteza
derivate din aceleași rotații ale roților nu sunt introduse simultan în filtru,
deoarece aceasta ar însemna folosirea de două ori a aceleiași informații.
Orientarea absolută a IMU nu este fuzionată: în pluginul Gazebo ea provine din
atitudinea exactă simulată și ar constitui o informație privilegiată. Filtrul
funcționează în modul plan \texttt{two\_d\_mode}, la frecvența nominală de
\SI{30}{Hz}. Rezultatul \texttt{/odometry/filtered} este singura stare folosită
de regulator. EKF-ul este și singurul proprietar al transformării dinamice
\texttt{odom} către \texttt{base\_footprint}.

Pentru evaluare se publică separat \texttt{/ground\_truth/odom}, calculată din
poziția și viteza corpului în lumea Gazebo. Această mărime nu este conectată la
regulator sau la EKF. Ea reprezintă un instrument de măsurare extern, folosit
numai pentru calcularea performanței reale și pentru analiza erorii de
localizare. Astfel, o derivă a localizării modifică reacția regulatorului și
poate deteriora mișcarea fizică, dar estimatorul nu își evaluează propria
estimare ca și cum aceasta ar fi adevărul.

\section{Fluxul informațional}

\begin{center}
\fbox{%
\begin{minipage}{0.86\textwidth}
\centering
Fișier CSV cu punctele căii\\
$\downarrow$\\
Manager comun de referință\\
$\downarrow$\\
Regulator selectat: PID lookahead sau PID cascadat\\
$\downarrow$\\
Comanda $(v_{\mathrm{cmd}},\omega_{\mathrm{cmd}})$ pe \texttt{/cmd\_vel}\\
$\downarrow$\\
Gazebo și modelul diferențial\\
$\downarrow$\\
Odometrie + IMU $\rightarrow$ EKF $\rightarrow$ \texttt{/odometry/filtered}\\
$\circlearrowleft$
\end{minipage}}
\end{center}

În paralel cu această buclă închisă, starea exactă Gazebo este transmisă numai
evaluatorului. Cele două ramuri au roluri diferite: starea filtrată închide
bucla de reacție, iar starea exactă măsoară rezultatul fizic al experimentului.

Bucla de control este declanșată de fiecare mesaj de odometrie filtrată. Pentru
două mesaje acceptate consecutiv, pasul discret este

\begin{equation}
  \Delta t_k=t_k-t_{k-1},
\end{equation}

unde $t_k$ este marcajul temporal ROS al estimării curente. Deoarece este
activat timpul simulat, regulatorul urmează timpul Gazebo, nu timpul de perete
al calculatorului gazdă. La primul eșantion se folosește perioada nominală
$1/\SI{30}{Hz}=\SI{0.0333}{s}$. Un marcaj repetat sau descrescător este ignorat,
iar pentru o întrerupere mai mare de \SI{0.2}{s} memoria celor două regulatoare
este resetată și se folosește o singură perioadă nominală. Astfel sunt evitate
o derivată nedefinită și integrarea unei erori pe un interval în care starea nu
a fost observată.

% =============================================================================
\chapter{Managerul comun de referință}
% =============================================================================

\section{Motivație}

O comparație corectă necesită ca PID, LQR și MPC să interpreteze identic calea și
erorile. Din acest motiv, prelucrarea geometrică a fost separată de regulator
într-o componentă fără dependențe ROS, denumită
\texttt{PathReferenceManager}. Componenta încarcă traseul, proiectează robotul
pe traseu, memorează progresul și furnizează aceleași mărimi de referință
oricărui regulator.

\section{Reprezentarea căii}

Fișierul de intrare este un CSV fără antet, cu exact două coloane finite
$(x_i,y_i)$. Sunt necesare minimum două puncte, iar punctele consecutive identice
sunt respinse deoarece ar genera segmente de lungime zero.

Pentru segmentul dintre punctele $\mathbf{p}_i$ și $\mathbf{p}_{i+1}$ se definesc:

\begin{align}
  \mathbf{d}_i &= \mathbf{p}_{i+1}-\mathbf{p}_i, \\
  \ell_i &= \lVert\mathbf{d}_i\rVert_2, \\
  \psi_i &= \operatorname{atan2}(d_{i,y},d_{i,x}).
\end{align}

Lungimile sunt cumulate pentru a obține coordonata curbilinie $s$ și lungimea
totală $S$.

\section{Proiecția robotului pe cale}

Pentru poziția robotului $\mathbf{p}$, proiecția neconstrânsă pe dreapta
segmentului este descrisă prin parametrul:

\begin{equation}
  t_i^*=\frac{(\mathbf{p}-\mathbf{p}_i)^T\mathbf{d}_i}
  {\mathbf{d}_i^T\mathbf{d}_i}.
\end{equation}

Deoarece traseul este format din segmente finite, valoarea este limitată:

\begin{equation}
  t_i=\sat_{[0,1]}(t_i^*), \qquad
  \mathbf{p}_{\mathrm{ref}}=\mathbf{p}_i+t_i\mathbf{d}_i.
\end{equation}

Parametrul $t_i$ precizează unde se află proiecția în interiorul segmentului:
$t_i=0$ la început, $t_i=1$ la sfârșit și valori intermediare între acestea.
Progresul continuu este:

\begin{equation}
  s=s_i+t_i\ell_i,
\end{equation}

unde $s_i$ este lungimea cumulată până la începutul segmentului.

Căutarea este efectuată numai într-o fereastră înaintea segmentului curent, iar
$s$ nu este lăsat să descrească. La distanțe egale se permite trecerea numai
între două segmente adiacente care împart același vârf. Aceste reguli împiedică
saltul pe o ramură greșită la intersecția traseului în formă de opt și împiedică
declararea imediată a finalului pe traseele închise.

\section{Erorile geometrice}

Vectorul tangent unitar al segmentului este
$\mathbf{t}=[\cos\psi_{\mathrm{ref}},\sin\psi_{\mathrm{ref}}]^T$. Eroarea
laterală semnată este:

\begin{equation}
  e_y=t_x(y-y_{\mathrm{ref}})-t_y(x-x_{\mathrm{ref}}).
\end{equation}

Prin această convenție, $e_y>0$ înseamnă că robotul se află în stânga sensului
de parcurgere al căii. Eroarea comună de orientare este:

\begin{equation}
  e_\psi=\wrap(\psi_{\mathrm{ref}}-\psi).
\end{equation}

\section{Curbura și viteza de referință}

Curbura este aproximată prin variația orientării tangente raportată la distanța
dintre centrele segmentelor. La colțurile poligonale, valoarea este limitată la
$|\kappa|\leq\kappa_{\max}$, unde profilurile curente folosesc
$\kappa_{\max}=\SI{5}{m^{-1}}$.

Politica de viteză comună este:

\begin{align}
  f_\kappa &= \frac{1}{1+k_\kappa|\kappa_{\mathrm{ref}}|}, \\
  f_g &= \sat_{[0,1]}\left(\frac{S-s}{d_g}\right), \\
  v_{\mathrm{ref}} &= v_0 f_\kappa f_g, \\
  \omega_{\mathrm{ref}} &= v_{\mathrm{ref}}\kappa_{\mathrm{ref}}.
\end{align}

În configurația curentă, $v_0=\SI{0.4}{m/s}$,
$k_\kappa=0{,}5$ și $d_g=\SI{0.5}{m}$. Factorul $f_g$ reduce liniar viteza în
apropierea capătului, iar $f_\kappa$ reduce viteza pe porțiunile cu curbură mare.

% =============================================================================
\chapter{Regulatoarele PID implementate}
% =============================================================================

\section{Forma discretă a regulatorului PID}

Pentru o eroare generică $e[k]$, implementarea folosește:

\begin{align}
  I[k] &= \sat_{[-I_{\max},I_{\max}]}
          \left(I[k-1]+e[k]\Delta t\right), \\
  D[k] &= \frac{e[k]-e[k-1]}{\Delta t}, \\
  u[k] &= K_p e[k]+K_i I[k]+K_d D[k].
\end{align}

După inițializare sau resetare, termenul derivativ este zero la primul eșantion,
evitând un impuls artificial. Integrala memorată este limitată pentru reducerea
fenomenului de windup. Stările PID sunt resetate după o întrerupere temporală
mare a odometriei sau după activarea protecției de timeout.

\section{Varianta preliminară: PID cu punct anticipat}

Prima variantă funcțională alege un punct aflat cu un număr fix de eșantioane
înaintea poziției proiectate. Un PID liniar folosește distanța până la acel punct,
iar un PID unghiular folosește eroarea dintre orientarea robotului și direcția
către punct. Viteza liniară este înmulțită cu
$\max(0,\cos e_{\mathrm{lookahead}})$ pentru a opri translația când robotul este
orientat departe de țintă.

Această variantă este păstrată ca reper preliminar și poate fi selectată prin
\texttt{controller\_mode: lookahead}. Principalul ei dezavantaj este dependența
de densitatea punctelor: cinci puncte pot însemna distanțe fizice diferite pe
fișiere CSV eșantionate diferit. În plus, regulatorul nu controlează direct
eroarea laterală folosită drept criteriu de performanță.

\section{PID cascadat bazat pe eroarea laterală}

Regulatorul principal implementat în această etapă este selectat prin
\texttt{controller\_mode: cascade}. El conține două bucle ierarhice.

\subsection{Bucla exterioară de poziție laterală}

PID-ul exterior primește eroarea laterală:

\begin{equation}
  u_y=K_{p,y}e_y+K_{i,y}\int e_y\,dt+K_{d,y}\frac{de_y}{dt}.
\end{equation}

Corecția de orientare are semnul opus deoarece o eroare pozitivă înseamnă că
robotul este în stânga căii și trebuie orientat spre dreapta:

\begin{align}
  \Delta\psi_c &= \sat_{[-\Delta\psi_{\max},\Delta\psi_{\max}]}(-u_y), \\
  \psi_d &= \wrap(\psi_{\mathrm{ref}}+\Delta\psi_c).
\end{align}

Limitarea lui $\Delta\psi_c$ previne solicitarea unei orientări nerealiste atunci
când abaterea laterală este mare.

\subsection{Bucla interioară de orientare}

Eroarea buclei interioare este:

\begin{equation}
  e_{\psi,c}=\wrap(\psi_d-\psi).
\end{equation}

Ieșirea specifică regulatorului este numai corecția vitezei unghiulare:

\begin{equation}
  \Delta\omega_{\mathrm{PID}}=
  K_{p,\psi}e_{\psi,c}+K_{i,\psi}\int e_{\psi,c}\,dt+
  K_{d,\psi}\frac{de_{\psi,c}}{dt}.
\end{equation}

Anticiparea curburii și saturația nu mai fac parte din clasa PID, ci dintr-o
politică de comandă reutilizabilă de PID, LQR și MPC. Această separare asigură
că fiecare regulator va produce aceeași mărime manipulată,
$\Delta\omega$, în aceleași condiții experimentale.

Un avantaj important al structurii este reacția imediată la o rotație externă.
Dacă robotul este rotit în timp ce $e_y=0$, bucla exterioară cere încă
$\Delta\psi_c=0$, dar $e_{\psi,c}$ se modifică instantaneu. Bucla interioară
produce imediat o comandă unghiulară, fără să aștepte apariția abaterii laterale.

\subsection{Politica de comandă comună}

În regim normal, viteza liniară este determinată exclusiv de geometria căii și
de distanța rămasă, independent de regulator:

\begin{align}
  v_{\mathrm{cmd}} &= \sat_{[0,v_{\max}]}(v_{\mathrm{ref}}), \\
  \omega_{\mathrm{ff}} &= \kappa_{\mathrm{ref}}v_{\mathrm{cmd}}, \\
  \omega_{\mathrm{cmd}} &=
  \sat_{[-\omega_{\max},\omega_{\max}]}
  (\omega_{\mathrm{ff}}+\Delta\omega).
\end{align}

Reducerea continuă a vitezei în funcție de eroarea PID a fost eliminată. În
forma anterioară, un regulator cu erori mai mari era ajutat implicit printr-o
viteză mai mică, ceea ce ar fi alterat comparația. Se păstrează doar o protecție
comună: translația este oprită dacă $|e_y|\geq\SI{0.75}{m}$ sau dacă
$|e_\psi|\geq\SI{1.2}{rad}$. Corecția unghiulară rămâne activă pentru
recuperare, iar anticiparea devine zero deoarece $v_{\mathrm{cmd}}=0$.

\subsection{Parametrii profilului reglat}

Reglarea a fost realizată secvențial, în ordinea impusă de structura cascadată.
Mai întâi a fost reglată bucla interioară de orientare, menținând neschimbată
bucla exterioară. Ulterior au fost modificate câștigurile buclei exterioare,
folosind valorile deja selectate pentru bucla interioară. Tabelul
~\ref{tab:cascade-gains} prezintă atât valorile inițiale, cât și profilul
selectat. Acesta este un reglaj experimental documentat, nu rezultatul unei
optimizări globale.

\begin{table}[htbp]
  \centering
  \caption{Câștigurile inițiale și valorile selectate după reglare}
  \label{tab:cascade-gains}
  \begin{tabular}{lrrl}
    \toprule
    Parametru & Inițial & Final & Rol \\
    \midrule
    $K_{p,y}$ & $1{,}2$ & $1{,}5$ & reacție proporțională la $e_y$ \\
    $K_{i,y}$ & $0$ & $0$ & acțiune integrală laterală \\
    $K_{d,y}$ & $0{,}15$ & $0{,}20$ & amortizarea variației lui $e_y$ \\
    $K_{p,\psi}$ & $2{,}5$ & $3{,}0$ & reacție proporțională la orientare \\
    $K_{i,\psi}$ & $0$ & $0$ & acțiune integrală unghiulară \\
    $K_{d,\psi}$ & $0{,}15$ & $0{,}20$ & amortizarea erorii unghiulare \\
    $\Delta\psi_{\max}$ & $0{,}7$ & $0{,}7$ & limita corecției [rad] \\
    $v_{\max}$ & $1{,}0$ & $1{,}0$ & limita liniară [m/s] \\
    $\omega_{\max}$ & $1{,}5$ & $1{,}5$ & limita unghiulară [rad/s] \\
    \bottomrule
  \end{tabular}
\end{table}

Deoarece $K_{i,y}=K_{i,\psi}=0$, profilul final se comportă efectiv ca două
regulatoare PD în cascadă, completate de anticiparea curburii și de legea
comună a comenzilor. Implementarea păstrează forma PID generală pentru a permite
introducerea ulterioară a acțiunii integrale, însă lucrarea trebuie să raporteze
valorile efectiv utilizate și să nu atribuie integralelor un efect inexistent.

% =============================================================================
\chapter{Considerații de implementare și reproductibilitate}
% =============================================================================

Detaliile software nu reprezintă obiectul principal al lucrării, însă două
decizii sunt necesare pentru validitatea comparației. În primul rând, calculul
proiecției, al erorilor, al curburii și al referințelor de viteză este comun
regulatoarelor PID, LQR și MPC. Tot într-un strat comun sunt aplicate viteza
liniară, anticiparea curburii, protecția la abateri mari și saturațiile; fiecare
regulator furnizează numai corecția unghiulară. În al doilea rând, câștigurile și condițiile
experimentale sunt păstrate în configurații independente de cod. Astfel,
schimbarea regulatorului nu modifică definiția problemei evaluate.

Fiecare probă pornește cu o simulare și un estimator reinițializate, utilizează
timpul simulat și înregistrează aceeași stare, aceleași referințe și aceleași
comenzi. Sunt limitate comenzile și corecția de orientare, iar datele temporale
invalide conduc la resetarea memoriei regulatorului. Aceste aspecte sunt
menționate pentru reproductibilitate; descrierea claselor, a funcțiilor și a
testelor software detaliate este păstrată în documentația tehnică a proiectului,
nu în corpul principal al tezei.

% =============================================================================
\chapter{Metodologia de evaluare}
% =============================================================================

\section{Traseele nominale}

Infrastructura de benchmark include cinci căi:

\begin{enumerate}
  \item dreaptă de \SI{5}{m}, pentru verificarea vitezei și a absenței
        oscilațiilor laterale;
  \item curbă sinusoidală, pentru schimbări continue ale curburii;
  \item traseu cu colț de aproximativ $90^\circ$, pentru regim tranzitoriu sever;
  \item cerc, pentru regim curbiliniu staționar și revenirea la punctul inițial;
  \item traseu în formă de opt, pentru schimbarea semnului curburii și tratarea
        corectă a auto-intersecției.
\end{enumerate}

Fiecare rulare pornește o instanță Gazebo nouă în mod headless. Astfel, starea
robotului, starea EKF și timpul simulat nu se transferă între experimente.

\section{Scenarii pentru separarea buclelor}

Traseele nominale combină simultan efectele erorii de orientare, erorii
laterale și curburii. Pentru reglare au fost introduse două scenarii auxiliare,
care nu sunt tratate drept hărți suplimentare de benchmark:

\begin{enumerate}
  \item o dreaptă paralelă cu axa inițială a robotului, deplasată lateral cu
        \SI{0.20}{m}, pentru excitarea dominantă a buclei exterioare;
  \item o dreaptă care pornește din poziția robotului la un unghi de
        $45^\circ$, pentru excitarea dominantă a buclei interioare.
\end{enumerate}

Această separare evită reglarea simultană a tuturor parametrilor pe un traseu
complex, caz în care nu ar fi clar ce buclă a produs o îmbunătățire sau o
degradare. Vârful inițial de orientare din al doilea scenariu este impus de
condiția inițială și nu poate fi redus de câștiguri; sunt relevante integrala
erorii, depășirea și timpul de așezare ulterior.

\section{Procedura de reglare secvențială}

Procedura aplicată a fost:

\begin{enumerate}
  \item rularea profilului inițial pe toate cele cinci trasee, fără modificări;
  \item creșterea câștigurilor buclei interioare de la
        $(K_{p,\psi},K_{d,\psi})=(2{,}5,0{,}15)$ la $(3{,}0,0{,}20)$ și
        evaluarea pe eroarea inițială de $45^\circ$;
  \item păstrarea buclei interioare selectate și creșterea buclei exterioare de
        la $(K_{p,y},K_{d,y})=(1{,}2,0{,}15)$ la $(1{,}5,0{,}20)$;
  \item evaluarea buclei exterioare pe abaterea laterală inițială de
        \SI{0.20}{m};
  \item verificarea profilului candidat pe toate traseele, pentru a evita
        suprareglarea pe un singur scenariu;
  \item promovarea câștigurilor numai după verificarea compromisului dintre
        eroare, durată, activitatea comenzilor și saturație;
  \item repetarea probei de perturbație cu profilul final.
\end{enumerate}

În timpul reglării nu au fost modificate simultan viteza de referință și
câștigurile PID. Pentru comparația finală, politica de viteză și limitele sunt
înghețate în componenta comună, izolând efectul legii de feedback.

\section{Indicatorii de performanță}

Pentru durata experimentului $T$ sunt calculați:

\begin{align}
  e_{y,\mathrm{RMS}} &=
  \sqrt{\frac{1}{T}\int_0^T e_y^2(t)\,dt}, \\
  e_{\psi,\mathrm{RMS}} &=
  \sqrt{\frac{1}{T}\int_0^T e_\psi^2(t)\,dt}, \\
  e_f &= \sqrt{(x_g-x(T))^2+(y_g-y(T))^2}, \\
  J_u &= \int_0^T\left[
    \left(\frac{v_{\mathrm{cmd}}}{v_{\max}}\right)^2+
    \left(\frac{\omega_{\mathrm{cmd}}}{\omega_{\max}}\right)^2
  \right]dt.
\end{align}

Ultimul indicator descrie activitatea comenzilor cinematice. El nu este
interpretat drept energie electrică sau efort mecanic al motoarelor, deoarece
interfața regulatorului furnizează referințe de viteză, nu cupluri.

În scenariile de reglare s-a folosit și integrala erorii absolute:

\begin{equation}
  \operatorname{IAE}(e)=\int_0^T |e(t)|\,dt.
\end{equation}

Timpul de așezare a fost declarat numai după zece eșantioane consecutive cu
$|e_y|\leq\SI{0.01}{m}$ și $|e_\psi|\leq\SI{0.03}{rad}$. Cerința de persistență
împiedică declararea eronată a așezării la o singură traversare a pragului.

Se mai înregistrează valorile maxime absolute ale erorilor și comenzilor,
durata, numărul de eșantioane și succesul atingerii destinației. Regulatorul se
oprește la o eroare estimată de cel mult \SI{0.08}{m}. Succesul fizic independent
se verifică din adevărul Gazebo cu pragul \SI{0.15}{m}, egal cu jumătate din
lungimea șasiului. Pragul mai larg nu ascunde eroarea: distanța finală exactă și
eroarea de localizare rămân raportate numeric pentru fiecare rulare.

Indicatorii principali de urmărire se calculează din poziția exactă Gazebo,
proiectată pe aceeași cale ca starea utilizată de regulator. Jurnalul intern al
regulatorului, bazat pe EKF, se păstrează pentru explicarea comenzilor, dar nu
este folosit drept arbitru final al performanței. Separat se calculează eroarea
de localizare

\begin{align}
  e_{p,\mathrm{loc}}(t) &=
  \sqrt{\left(\hat{x}(t)-x_{\mathrm{GT}}(t)\right)^2+
  \left(\hat{y}(t)-y_{\mathrm{GT}}(t)\right)^2}, \\
  e_{\psi,\mathrm{loc}}(t) &=
  \operatorname{wrap}\left(\hat{\psi}(t)-\psi_{\mathrm{GT}}(t)\right).
\end{align}

Fluxurile au marcaje temporale proprii. Pentru calculul erorii de localizare,
poziția exactă este interpolată la marcajul temporal al estimării EKF. Fără
această aliniere, întârzierea normală dintre două mesaje ar apărea fals drept o
eroare de poziție aproximativ egală cu produsul dintre viteză și întârziere.

\section{Proba de perturbație}

Robustețea a fost evaluată printr-o eroare temporară introdusă în calea de
comandă, după regulator. Între secundele 5 și 6 s-a adăugat o polarizare
$\Delta\omega=+\SI{0.6}{rad/s}$ comenzii unghiulare, fără modificarea comenzii
liniare. Regulatorul nu primește informația că perturbația este activă; el
reacționează numai după ce abaterea devine observabilă în starea estimată.

Această alegere modelează o eroare temporară de acționare, calibrare sau
transmitere a comenzii și evită modificarea contactului roată--sol. Pentru toate
regulatoarele viitoare trebuie păstrate aceeași amplitudine, același interval și
aceleași limite de comandă. Revenirea strictă este declarată după zece
eșantioane consecutive cu $|e_y|\leq\SI{0.005}{m}$ și
$|e_\psi|\leq\SI{0.01}{rad}$, măsurate după sfârșitul perturbației.

Datele complete sunt păstrate pentru reprezentarea în MATLAB a căii reale,
erorilor, referințelor și comenzilor. Lista celor 26 de coloane și descrierea
funcțiilor software apar în documentația tehnică separată; în teză sunt reținute
numai mărimile necesare definirii și reproducerii indicatorilor.

% =============================================================================
\chapter{Rezultatele reglării și validării PID}
% =============================================================================

Tabelele de selectare a câștigurilor au fost obținute în etapa de punere în
funcțiune, înainte de separarea completă a politicii de viteză. Ele documentează
alegerea valorilor PID, dar nu sunt folosite drept comparație finală cu LQR și
MPC. După introducerea politicii comune, validarea nominală și proba de
perturbație au fost repetate cu evaluare din \texttt{/ground\_truth/odom}.
Rezultatele de mai jos provin din câte o rulare de verificare; matricea formală
va include repetări și indicatori de dispersie.

\section{Selectarea buclei interioare}

Tabelul~\ref{tab:inner-tuning} compară profilul inițial cu modificarea exclusivă
a buclei de orientare. Creșterea câștigurilor a redus ambele integrale ale
erorii și a scurtat timpul de așezare cu aproximativ \SI{0.29}{s}. A apărut un
singur eșantion suplimentar la limita vitezei unghiulare, fără creșterea
depășirii sau a activității totale a comenzilor. Din acest motiv au fost selectate
$K_{p,\psi}=3{,}0$ și $K_{d,\psi}=0{,}20$.

\begin{table}[htbp]
  \centering
  \caption{Reglarea buclei interioare pe eroarea inițială de $45^\circ$}
  \label{tab:inner-tuning}
  \begin{tabular}{lrrrr}
    \toprule
    Profil & IAE $e_\psi$ [rad s] & IAE $e_y$ [m s] & $t_s$ [s] & eșant. sat. \\
    \midrule
    Inițial & 0,5911 & 0,2117 & 4,092 & 4 \\
    Buclă interioară selectată & 0,5434 & 0,1943 & 3,800 & 5 \\
    \bottomrule
  \end{tabular}
\end{table}

\section{Selectarea buclei exterioare}

După fixarea buclei interioare, câștigurile laterale au fost evaluate pe
abaterea inițială de \SI{0.20}{m}. Rezultatele din
Tabelul~\ref{tab:outer-tuning} arată reducerea IAE laterale cu aproximativ
$17\%$ și a timpului de așezare de la \SI{7.0}{s} la \SI{5.6}{s}. IAE de
orientare a rămas practic neschimbată. Creșterea comenzii unghiulare maxime de
la \SI{0.72}{rad/s} la \SI{0.90}{rad/s} a rămas sub limita de
\SI{1.5}{rad/s}. Au fost astfel selectate $K_{p,y}=1{,}5$ și
$K_{d,y}=0{,}20$.

\begin{table}[htbp]
  \centering
  \caption{Reglarea buclei exterioare pe abaterea laterală de \SI{0.20}{m}}
  \label{tab:outer-tuning}
  \begin{tabular}{lrrrr}
    \toprule
    Câștiguri $(K_{p,y},K_{d,y})$ & IAE $e_y$ [m s] & IAE $e_\psi$ [rad s]
      & $t_s$ [s] & $|\omega|_{\max}$ [rad/s] \\
    \midrule
    $(1{,}2,\ 0{,}15)$ & 0,5525 & 0,5811 & 7,0 & 0,720 \\
    $(1{,}5,\ 0{,}20)$ & 0,4568 & 0,5814 & 5,6 & 0,900 \\
    \bottomrule
  \end{tabular}
\end{table}

\section{Validarea profilului final pe traseele nominale}

Tabelul~\ref{tab:cascade-results} prezintă rularea de validare de după separarea
politicii comune. Toate traseele au fost finalizate. Eroarea estimată finală apropiată de
\SI{0.08}{m} este determinată în principal de toleranța de oprire
$\SI{0.08}{m}$ și nu trebuie confundată cu eroarea laterală staționară.

\begin{table}[htbp]
  \centering
  \small
  \caption{Validarea PID după introducerea politicii comune; erori fizice Gazebo}
  \label{tab:cascade-results}
  \begin{tabular}{lrrrrrr}
    \toprule
    Traseu & Succes & $T$ [s] & RMS $e_y$ [m] & max $|e_y|$ [m]
      & RMS $e_\psi$ [rad] & $J_u$ \\
    \midrule
    Dreaptă & da & 13,300 & 0,00176 & 0,00251 & 0,00052 & 1,961 \\
    Curbă sinusoidală & da & 19,670 & 0,02169 & 0,10750 & 0,18127 & 4,446 \\
    Colț & da & 18,200 & 0,08594 & 0,25460 & 0,40742 & 5,757 \\
    Cerc & da & 24,798 & 0,00664 & 0,01802 & 0,06785 & 3,202 \\
    Opt & da & 40,300 & 0,03342 & 0,09056 & 0,11372 & 6,867 \\
    \bottomrule
  \end{tabular}
\end{table}

Rezultatele nu se compară numeric cu vechea arhitectură, deoarece aceasta reducea
viteza în funcție de eroarea PID. Colțul rămâne cazul cu cea mai mare eroare deoarece
referința poligonală impune o variație instantanee de aproximativ $90^\circ$,
imposibil de urmărit instantaneu de un sistem cu viteză și accelerație limitate.
Protecția comună a oprit translația timp de 29 de eșantioane în vecinătatea
colțului, fără a împiedica finalizarea. Pe celelalte trasee nu a fost activată.

\section{Rejecția perturbației de comandă}

Tabelul~\ref{tab:disturbance-results} prezintă repetarea probei după separarea
politicii comune de comandă.

\begin{table}[htbp]
  \centering
  \caption{Răspunsul la polarizarea unghiulară de \SI{0.6}{rad/s}, aplicată
  timp de \SI{1}{s}}
  \label{tab:disturbance-results}
  \begin{tabular}{lrrrr}
    \toprule
    Profil & max $|e_y|$ [m] & max $|e_\psi|$ [rad] & revenire [s]
      & $|e_y|$ final [m] \\
    \midrule
    PID cu politică comună & 0,05729 & 0,13870 & 4,361 & 0,00312 \\
    \bottomrule
  \end{tabular}
\end{table}

Rezultatul confirmă funcționarea ambelor niveluri ale cascadei. Polarizarea
produce inițial o eroare de orientare, la care bucla interioară reacționează
direct. Deplasarea laterală rezultată modifică apoi orientarea dorită prin bucla
exterioară. Totuși, valorile provin dintr-o singură rulare și au rol
de validare a reglajului; concluziile comparative finale necesită repetări și
raportarea dispersiei.

\section{Necesitatea acțiunii integrale}

Deoarece profilul selectat are $K_{i,y}=K_{i,\psi}=0$, verificarea relevantă a
folosit polarizări unghiulare constante și temporare în calea de comandă. O
probă anterioară pe o pantă de $15^\circ$ a fost exclusă din metodologia finală:
odometria roților integrează distanța parcursă pe suprafața înclinată, în timp
ce referința plană și adevărul Gazebo folosesc proiecția orizontală. În acea
configurație, regulatorul putea declara atingerea capătului înainte ca robotul
să ajungă fizic la punctul final, deci proba amesteca eroarea de localizare cu
performanța regulatorului.

Verificarea păstrată a aplicat pe dreaptă o polarizare unghiulară constantă
$b=\SI{0.08}{rad/s}$ între secundele 3 și 10. Pentru regimul staționar pe o
dreaptă, neglijând termenii derivați și considerând orientarea corpului paralelă
cu traseul, rezultă:

\begin{align}
  e_{\psi,c} &\simeq -K_{p,y}e_y, \\
  \omega_{\mathrm{aplicată}} &\simeq
    -K_{p,\psi}K_{p,y}e_y+b.
\end{align}

Condiția $\omega_{\mathrm{aplicată}}=0$ conduce la abaterea laterală necesară
pentru generarea corecției proporționale:

\begin{equation}
  e_{y,\mathrm{ss}}=\frac{b}{K_{p,\psi}K_{p,y}}
  =\frac{0{,}08}{3{,}0\cdot1{,}5}=\SI{0.01778}{m}.
\end{equation}

În ultimele două secunde ale polarizării s-a măsurat o medie
$\overline{e_y}=\SI{0.01731}{m}$, cu abaterea standard
\SI{0.00017}{m}. Comanda nominală unghiulară medie a fost
$-\SI{0.08047}{rad/s}$, anulând aproape complet polarizarea. Concordanța cu
predicția confirmă că, în absența acțiunii integrale, o perturbație constantă
este compensată printr-o abatere staționară mică, dar nenulă.

Pentru a verifica dacă rezultatul se menține la o polarizare mai mare, a fost
aplicat $b=\SI{0.20}{rad/s}$ timp de \SI{15}{s} pe o dreaptă de \SI{10}{m}.
Au fost comparate valorile $K_{i,\psi}\in\{0;0{,}2;0{,}4;0{,}6\}$, cu aceeași
limită a stării integrale $I_{\max,\psi}=0{,}5$. Rezultatele sunt prezentate în
Tabelul~\ref{tab:integral-constant-bias}.

\begin{table}[htbp]
  \centering
  \caption{Compromisul acțiunii integrale la polarizarea constantă de
  \SI{0.20}{rad/s}}
  \label{tab:integral-constant-bias}
  \begin{tabular}{rrrrr}
    \toprule
    $K_{i,\psi}$ & IAE activă [m s] & $\overline{e_y}$ final activ [m]
      & IAE după eliminare [m s] & revenire [s] \\
    \midrule
    0,0 & 0,5820 & 0,04445 & 0,0849 & 3,801 \\
    0,2 & 0,3721 & 0,02286 & 0,1418 & nerevenit \\
    0,4 & 0,2695 & 0,00976 & 0,1705 & nerevenit \\
    0,6 & 0,2032 & 0,00464 & 0,1557 & nerevenit \\
    \bottomrule
  \end{tabular}
\end{table}

Acțiunea integrală îmbunătățește fără echivoc rejecția cât timp defectul rămâne
activ. Pentru $K_{i,\psi}=0{,}6$, eroarea medie din ultimele două secunde scade
de la \SI{44.45}{mm} la \SI{4.64}{mm}, iar IAE activă scade cu aproximativ
$65\%$. După dispariția polarizării, starea integrală memorată continuă însă să
producă o comandă în sens opus și crește eroarea de revenire.

Limitarea integralei nu elimină acest compromis. Pentru $K_{i,\psi}=0{,}2$,
starea estimată a ajuns la limita $-0{,}5$, dar pentru valorile $0{,}4$ și
$0{,}6$ stările finale estimate au fost aproximativ $-0{,}417$, respectiv
$-0{,}311$, deci nu erau saturate. Comanda unghiulară nu a atins limita
actuatorului. În ultimele două cazuri nu este vorba despre windup produs de
saturație, ci despre memoria necesară compensării defectului constant; după
eliminarea defectului, aceasta trebuie descărcată printr-o eroare de semn opus.

Profilul $K_{i,\psi}=0{,}6$ a fost verificat și pe trasee nominale. RMS-ul
erorii laterale a scăzut de la \SI{6.82}{mm} la \SI{2.16}{mm} pe cerc și de la
\SI{6.16}{mm} la \SI{3.75}{mm} pe traseul în formă de opt, cu o creștere a
activității comenzilor sub aproximativ $1\%$. În schimb, la perturbația principală de
$\SI{0.6}{rad/s}$ aplicată timp de o secundă, profilul fără integrală a revenit
în \SI{4.404}{s} și a încheiat cu \SI{0.11}{mm} abatere laterală, în timp ce
profilul cu $K_{i,\psi}=0{,}6$ nu a îndeplinit criteriul de revenire și a
încheiat cu \SI{8.29}{mm}.

Prin urmare, niciuna dintre valorile integrale testate nu domină profilul fără
integrală pentru toate criteriile. Dacă defectul constant este obiectivul
prioritar, $K_{i,\psi}=0{,}6$ oferă robustețe staționară superioară. Pentru
metodologia curentă, în care perturbația principală este temporară și timpul de
revenire este un criteriu explicit, a fost păstrat $K_{i,\psi}=0$. Această
decizie este un compromis experimental, nu afirmația că acțiunea integrală ar
fi inutilă în general.

\section{Verificarea funcțională}

Înaintea experimentelor, definițiile geometrice, semnele cascadei, reacția la o
rotație fără abatere laterală, feedforward-ul de curbură, saturațiile și
frontierele temporale ale perturbației au fost verificate automat. Cele 14 ținte
CTest ale celor două pachete se încheie fără erori sau eșecuri. Instrumentul
ROS omite 18 verificări statice cppcheck deoarece versiunea instalată este
marcată drept lentă; testele geometrice și de control sunt executate. Acest rezultat
este menționat ca verificare a corectitudinii experimentului, fără a detalia în
teză fiecare funcție sau caz software.

% =============================================================================
\chapter{Stadiul proiectului și etapele următoare}
% =============================================================================

Până la această etapă au fost realizate:

\begin{itemize}
  \item corectarea geometriei roților și a responsabilităților TF;
  \item configurarea estimării plane prin EKF și folosirea timpului simulat;
  \item protecții la date invalide, întreruperi de odometrie și oprire;
  \item păstrarea PID-ului lookahead ca reper preliminar;
  \item profiluri baseline, fast și robust pentru această variantă;
  \item trasee dreaptă, sinusoidală, colț, cerc și opt;
  \item script de benchmark și indicatori calculați automat;
  \item manager comun de referință cu proiecție și progres continuu;
  \item politică comună de viteză, anticipare, siguranță și saturație pentru
        PID, LQR și MPC;
  \item implementarea și testarea PID-ului cascadat;
  \item reglarea secvențială și înghețarea câștigurilor profilului PID;
  \item validarea pe cinci trasee și două condiții inițiale izolate;
  \item implementarea și repetarea unei perturbații de comandă reproductibile;
  \item jurnalizarea variabilelor necesare analizei în MATLAB.
\end{itemize}

Etapele recomandate în continuare sunt:

\begin{enumerate}
  \item repetarea fiecărui experiment PID și raportarea mediei, abaterii standard și
        eventual a intervalelor de încredere;
  \item realizarea scripturilor MATLAB pentru suprapunerea căii reale peste
        referință și pentru graficele erorilor și comenzilor;
  \item derivarea și validarea modelului analitic liniarizat în coordonate de
        eroare față de cale;
  \item implementarea LQR folosind acest model și același manager de referință;
  \item implementarea MPC cu aceleași limite și criterii;
  \item evaluarea finală comparativă PID--LQR--MPC.
\end{enumerate}

Limitările curente sunt importante pentru interpretare. Calea este poligonală,
iar curbura este o aproximație discretă. Reglajul cascadei este experimental,
nu rezultatul unei optimizări globale. Rezultatele din acest document provin
din câte o rulare și nu oferă încă o bază statistică. În plus, problema curentă
este de urmărire geometrică a căii, fără o referință temporală strictă.
Identificarea experimentală a sistemului a fost eliminată din domeniul lucrării.
Pluginul Gazebo aplică o servocomandă de viteză cu o limită de accelerație
explicit cunoscută, iar proiectul nu urmărește transferul pe un robot fizic;
LQR și MPC vor folosi modelul cinematic analitic. LQR și MPC nu fac încă parte
din implementarea descrisă aici.

% =============================================================================
\appendix
\chapter{Fișiere relevante}
% =============================================================================

\begin{longtable}{p{0.40\textwidth}p{0.52\textwidth}}
  \toprule
  Fișier & Responsabilitate \\
  \midrule
  \endhead
  \path{src/my_robot_controller/src/path_reference_manager.cpp} &
  geometria, proiecția, progresul, curbura și referințele; \\
  \path{src/my_robot_controller/src/pid_controller.cpp} &
  PID discret generic și anti-windup; \\
  \path{src/my_robot_controller/src/cascaded_pid_controller.cpp} &
  legea de feedback PID cascadată; \\
  \path{src/my_robot_controller/src/motion_command_policy.cpp} &
  viteza, anticiparea curburii, protecția și limitele comune regulatoarelor; \\
  \path{src/my_robot_controller/src/pid_node.cpp} &
  integrarea ROS, selectarea modului și jurnalizarea; \\
  \path{src/my_robot_controller/config/pid_cascade.yaml} &
  câștigurile și limitele profilului cascadat; \\
  \path{scripts/run_pid_benchmarks.sh} &
  automatizarea experimentelor și calculul indicatorilor; \\
  \path{results/cascade_pid_pre_tuning_20260720/summary.csv} &
  rezultatele profilului inițial, înainte de reglare; \\
  \path{results/cascade_pid_final_20260720/summary.csv} &
  validarea profilului final pe trasee și condiții inițiale; \\
  \path{results/cascade_command_disturbance_final_20260720/recovery_summary.txt} &
  rezultatul final al probei de perturbație; \\
  \path{results/shared_motion_policy_nominal_20260807/ground_truth_summary.csv} &
  verificarea nominală după separarea politicii comune; \\
  \path{results/shared_motion_policy_disturbance_20260807/recovery_summary.txt} &
  proba de perturbație după separarea politicii comune; \\
  \path{results/cascade_constant_bias_diagnostic_20260720/steady_state_summary.txt} &
  verificarea regimului staționar fără acțiune integrală; \\
  \path{results/integral_gain_study_20260720.csv} &
  comparația câștigurilor integrale la defect constant și tranzitoriu. \\
  \bottomrule
\end{longtable}

\end{document}
